% Report acronyms. Mathematical notation is defined in notation.tex.
% \syma uses the nomencl A (abbreviation) group defined there.
\cleardoublepage
\renewcommand{\nomname}{List of Abbreviations}
\markboth{\nomname}{\nomname}
\renewcommand{\nompreamble}{\addcontentsline{toc}{chapter}{\nomname}%
\setlength{\nomitemsep}{0.2\baselineskip}
}

\syma{AIV}{Assembly, integration and verification}
\syma{API}{Application programming interface}
\syma{BM25}{Lexical ranking method for text retrieval}
\syma{CAD}{Computer-aided design}
\syma{CDR}{Critical design review}
\syma{cFS}{core Flight System}
\syma{ConOps}{Concept of operations}
\syma{COTS}{Commercial off-the-shelf}
\syma{CPU}{Central processing unit}
\syma{CSRM}{CubeSat System Reference Model}
\syma{DC}{Direct current}
\syma{DFN}{Doyle-Fuller-Newman battery model}
\syma{DSM}{Design structure matrix}
\syma{DUT}{Device under test}
\syma{EMC}{Electromagnetic compatibility}
\syma{EPS}{Electrical power system}
\syma{FMEA}{Failure modes and effects analysis}
\syma{FMECA}{Failure modes, effects and criticality analysis}
\syma{FMI}{Functional Mock-up Interface}
\syma{FMU}{Functional Mock-up Unit}
\syma{FN}{False negative}
\syma{FP}{False positive}
\syma{FPP}{F Prime Prime (F Prime modeling language)}
\syma{FR}{Framework requirement (system requirement for the EPS engineering framework)}
\syma{FRR}{Flight readiness review}
\syma{FSW}{Flight software}
\syma{FW}{Framework (action-label prefix)}
\syma{GNSS}{Global navigation satellite system}
\syma{GPS}{Global Positioning System}
\syma{GPU}{Graphics processing unit}
\syma{HDDL}{Hierarchical domain definition language}
\syma{HIL}{Hardware in the loop}
\syma{HTN}{Hierarchical task network}
\syma{ID}{Identifier}
\syma{IF}{Interface}
\syma{IMU}{Inertial measurement unit}
\syma{INCOSE}{International Council on Systems Engineering}
\syma{KiBaM}{Kinetic battery model}
\syma{LED}{Light-emitting diode}
\syma{MBSE}{Model-based systems engineering}
\syma{MCR}{Mission concept review}
\syma{MDR}{Mission definition review}
\syma{MIL}{Model in the loop}
\syma{MOE}{Measure of effectiveness}
\syma{MOP}{Measure of performance}
\syma{MPPT}{Maximum power point tracking}
\syma{MR}{Mission requirement}
\syma{NASA}{National Aeronautics and Space Administration}
\syma{NE}{Non-EPS engineer (stakeholder role)}
\syma{OBC}{On-board computer}
\syma{OMG}{Object Management Group}
\syma{OOSEM}{Object-Oriented Systems Engineering Method}
\syma{PCB}{Printed circuit board}
\syma{PCI}{Peripheral Component Interconnect}
\syma{PDR}{Preliminary design review}
\syma{PIL}{Processor in the loop}
\syma{PM}{Project manager (stakeholder role)}
\syma{RAM}{Reliability, availability and maintainability}
\syma{REST}{Representational state transfer}
\syma{RMSE}{Root mean square error}
\syma{SaaS}{Software as a service}
\syma{SC}{Scenario}
\syma{SDR}{System definition review}
\syma{SE}{Systems engineering. Systems engineer when naming a stakeholder role.}
\syma{SIL}{Software in the loop}
\syma{SIR}{System integration review}
\syma{SOH}{State of health}
\syma{SPM}{Single particle battery model}
\syma{SRR}{System requirements review}
\syma{SysML}{Systems Modeling Language}
\syma{TBR}{To be resolved}
\syma{TP}{True positive}
\syma{UC}{Use case}
\syma{USD}{Universal Scene Description (scene representation)}
\syma{V\&V}{Verification and validation}
\syma{YAML}{YAML Ain't Markup Language}

\printnomenclature[20mm]
